
\subsubsection{Spurious Correlation Detection}

Geirhos et al. (2020) surveyed spurious correlation mechanisms, showing that networks preferentially learn texture over shape in ImageNet and background over foreground in Waterbirds. Sagawa et al. (2020) introduced worst-group accuracy as a metric that reveals spurious reliance: on Waterbirds, Empirical Risk Minimization (ERM) achieves 97\% average accuracy but 72.6\% worst-group accuracy. Group DRO improves worst-group accuracy to 91.4\% by reweighting loss toward minority groups. These methods evaluate at convergence only and do not explain temporal dynamics or architectural effects.

\subsubsection{Temporal Learning Dynamics}

Toneva et al. (2019) introduced ''example forgetting'' to characterize temporal learning: unforgettable examples (learned early, never misclassified again) tend to be simple or spurious, while forgettable examples contain complex core features. On CIFAR-10, 30\% of examples are unforgettable. Arpit et al. (2017) showed that networks memorize training examples in order of increasing difficulty, with easy (spurious) examples memorized first. These works establish that temporal dynamics exist but do not compare architectural components or provide gradient-level explanations.

\subsubsection{Batch Normalization}

Batch Normalization \citep{Ioffe2015Batch} normalizes activations using batch statistics (mean and variance computed over the batch dimension). Santurkar et al. (2019) demonstrated that Batch Normalization smooths the loss landscape, enabling faster convergence, rather than reducing internal covariate shift. Shen et al. (2021) observed that Batch Normalization can harm worst-group accuracy, hypothesizing that batch statistics encode spurious batch-level correlations, but provided no controlled comparison or mechanistic explanation. Layer Normalization \citep{Ba2016Layer} normalizes per-instance over feature dimensions rather than per-batch over samples. While standard in Transformers, its effect on spurious correlation learning in convolutional networks has not been systematically studied.

This work provides the first controlled comparison of Batch Normalization and Layer Normalization on spurious correlation tasks, with gradient-level mechanistic analysis and accuracy-matched temporal comparison.
